\chapter{Das Unternehmen Erpolino}
\label{kap:erpolino}

Die vorangegangenen Kapitel haben das Zuordnungsproblem in allgemeiner Form
entwickelt. Damit aus dem Modell eine rechenbare Instanz wird, sind konkrete
Zahlen nötig: Wer arbeitet im Betrieb, welche Tätigkeiten darf diese Person
ausführen, wie schnell erledigt sie diese, und welche Arbeitsgänge stehen
überhaupt an. Dieses Kapitel beschreibt das Unternehmen, aus dem diese Daten
stammen, und legt offen, wie die Zahlen der Untersuchung zustande
kommen.\footnote{Die Erpolino Metallverarbeitung GmbH ist ein frei erfundenes
Unternehmen. Sämtliche in dieser Arbeit genannten Angaben zu Geschichte,
Größe, Kundenstruktur, Personal und Auftragsbestand sowie alle Zahlenwerte der
Qualifikationsmatrix und des Fertigungsloses sind konstruiert und dienen
ausschließlich dazu, das Verfahren an einer in sich stimmigen Instanz
vorzuführen. Ein realer Betrieb dieses oder ähnlichen Namens ist nicht
gemeint, und die Ergebnisse lassen keinen Rückschluss auf reale Betriebsdaten
zu.}

\section{Unternehmen und Fertigung}
\label{sec:unternehmen}

Die Erpolino Metallverarbeitung GmbH mit Sitz in Reutlingen wird 1978 als
Dreherei mit vier Beschäftigten gegründet. Der Standort liegt in der Region
Neckar-Alb, einem von mittelständischer Zulieferindustrie geprägten Teil
Baden-Württembergs, dessen Nähe zu den Ballungsräumen Stuttgart und Tübingen
den Kundenkreis des Unternehmens bis heute prägt. Seit 2004 führt die zweite
Generation der Gründerfamilie das Unternehmen; die Geschäftsanteile befinden
sich vollständig in Familienbesitz. Diese Eigentümerstruktur ist für die
Fragestellung der vorliegenden Arbeit insofern relevant, als Investitions- und
Organisationsentscheidungen ohne externe Kapitalgeber getroffen werden und
Veränderungen in der Planung entsprechend selten aus formalen Vorgaben,
sondern aus praktischem Bedarf entstehen.

Gegenwärtig beschäftigt das Unternehmen rund 120 Mitarbeitende, von denen etwa
45 in der Fertigung tätig sind. Die übrigen Stellen verteilen sich auf
Konstruktion, Vertrieb, Einkauf, Qualitätsmanagement, Verwaltung und
Instandhaltung. Das Unternehmen betreibt keine eigene Produktentwicklung und
vertreibt keine eigenen Erzeugnisse, sondern arbeitet ausschließlich in der
\Index{Lohnfertigung}: Es fertigt Präzisionsteile nach Zeichnungen und
Spezifikationen seiner Auftraggeber. Der wirtschaftliche Kern des
Geschäftsmodells liegt damit nicht im Produkt, sondern in der verfügbaren
Fertigungskapazität und in der Fähigkeit, diese termingerecht bereitzustellen.

Die Kunden stammen im Wesentlichen aus drei Branchen. Der \Index{Maschinenbau}
bezieht Gehäuse-, Trag- und Rahmenteile in kleinen Stückzahlen, häufig als
Ersatz- oder Vorserienteile. Die \Index{Antriebstechnik} fragt Wellen,
Buchsen, Lagerböcke und Getriebedeckel nach, meist in mittleren Losgrößen und
mit engen Form- und Lagetoleranzen. Die \Index{Medizintechnik} schließlich
verlangt kleine Stückzahlen bei hohen Anforderungen an Oberflächengüte und
Dokumentation. Keiner dieser Kunden vergibt Großserien; die typische
\Index{Losgröße} liegt zwischen einem Dutzend und einigen hundert Teilen. Für
die Fertigung bedeutet dies einen häufigen Wechsel zwischen Aufträgen und
einen entsprechend hohen Anteil an Rüst- und Einrichtvorgängen.

Die Fertigungstiefe umfasst sechs Bereiche: Drehen, Fräsen, Schweißen,
Schleifen, CNC-Programmierung und Qualitätsprüfung. Der Maschinenpark besteht
überwiegend aus CNC-Dreh- und Bearbeitungszentren, ergänzt um konventionelle
Maschinen für Einzelteile und Nacharbeit. Wärmebehandlung und
Oberflächenveredelung werden zugekauft und sind nicht Gegenstand dieser
Arbeit. Die sechs Bereiche sind keine räumlich getrennten Abteilungen mit
festem Personal, sondern Tätigkeitsfelder, zwischen denen die Beschäftigten je
nach Auftragslage wechseln. Genau diese Flexibilität ist der Ausgangspunkt der
Untersuchung: Sie eröffnet Spielräume in der Zuteilung, die bislang nicht
systematisch genutzt werden.

\section{Qualifikationen der Mitarbeitenden}
\label{sec:qualifikationen}

Ob eine Person einen Arbeitsgang übernehmen kann, entscheidet sich an ihrer
Qualifikation. Das Unternehmen unterscheidet die \ZahlQualifikationen{}
Qualifikationen Drehen, Fräsen, Schweißen, Schleifen, CNC-Programmierung und
Qualitätsprüfung. Sie werden über Ausbildung, innerbetriebliche Einweisung und
in Teilen über externe Zertifikate erworben; das Schweißen und die
Qualitätsprüfung unterliegen dabei formalen Nachweispflichten gegenüber den
Kunden der Medizintechnik. Niemand in der Fertigung beherrscht alle sechs
Felder, und niemand beherrscht nur eines: Die Beschäftigten verfügen in der
Regel über zwei oder drei Qualifikationen.

Erfasst wird dieser Sachverhalt in einer \Index{Qualifikationsmatrix}. Sie ordnet
jedem Paar aus Mitarbeiter und Qualifikation entweder einen Zahlenwert oder die
Angabe zu, dass die Qualifikation nicht vorliegt. Der Zahlenwert ist die
\Index{Effizienz} der Person in diesem Feld. Sie ist als Verhältniszahl definiert:
Der Wert $1{,}00$ entspricht der \Index{Normalleistung}, also derjenigen Leistung, die
der Arbeitsplanung als Vorgabewert zugrunde liegt. Ein Wert über $1{,}00$
bedeutet, dass die Person schneller arbeitet als die Vorgabe, ein Wert
darunter, dass sie länger braucht. Die Bearbeitungsdauer eines Arbeitsgangs
ergibt sich daraus als Quotient aus der Grundzeit des Arbeitsgangs und der
Effizienz der ausführenden Person. Eine Grundzeit von 120 Minuten führt bei
einer Effizienz von $1{,}20$ folglich auf 100 Minuten, bei einer Effizienz von
$0{,}90$ dagegen auf rund 133 Minuten.

Die im Unternehmen auftretenden Werte liegen zwischen $0{,}85$ und $1{,}40$.
Diese Spanne ist erklärungsbedürftig, denn sie ist erheblich: Zwischen der
langsamsten und der schnellsten Bearbeitung desselben Arbeitsgangs liegt ein
Faktor von mehr als eineinhalb. Sie erklärt sich aus der Zusammensetzung der
Belegschaft. Wer eine Tätigkeit seit Jahren täglich ausübt, erreicht Werte im
oberen Bereich; wer die Qualifikation zwar besitzt, sie aber nur gelegentlich
einsetzt, liegt unter der Normalleistung. Die Effizienz misst damit keine
persönliche Leistungsbereitschaft, sondern Routine im jeweiligen Feld. Sie ist
zudem qualifikationsbezogen und nicht personenbezogen: Dieselbe Person kann in
einem Feld über und in einem anderen unter der Normalleistung liegen.

Fehlt die Qualifikation, so ist die Zuordnung nicht etwa langsam, sondern
unzulässig. Für diesen Fall existiert kein Effizienzwert, und in der Matrix
steht ein Strich. Diese Unterscheidung ist für die Modellierung wesentlich: Ein
fehlender Eintrag darf nicht durch einen sehr kleinen Effizienzwert ersetzt
werden, weil das Verfahren eine solche Zuordnung sonst zwar teuer, aber
grundsätzlich zulässig behandeln würde. Wie diese Schranke im linearen Modell
umgesetzt wird, ist in Kapitel~3 dargelegt.

\begin{landscape}
\begin{table}
  \centering
  % Automatisch erzeugt von Code/tabellen.py -- nicht von Hand ändern.
% Quelle: Code/daten.py und Ergebnisse/*.json

\begin{tabular}{@{}lS[table-format=1.2]S[table-format=1.2]S[table-format=1.2]S[table-format=1.2]S[table-format=1.2]S[table-format=1.2]S[table-format=1.0]@{}}
\toprule
Mitarbeiter & {\rotatebox{45}{Drehen}} & {\rotatebox{45}{Fräsen}} & {\rotatebox{45}{Schweißen}} & {\rotatebox{45}{Schleifen}} & {\rotatebox{45}{CNC-Programmierung}} & {\rotatebox{45}{Qualitätsprüfung}} & {\rotatebox{45}{Qualifikationen}} \\
\midrule
  Bauer & 1.25 & 1.05 & {--} & 0.95 & {--} & {--} & 3 \\
  Cakir & 0.90 & 1.30 & {--} & {--} & 1.10 & {--} & 3 \\
  Dobrev & {--} & {--} & 1.35 & 1.00 & {--} & {--} & 2 \\
  Engelhardt & {--} & 1.00 & {--} & {--} & 1.40 & {--} & 2 \\
  Fischer & 1.10 & {--} & 0.95 & {--} & {--} & 1.15 & 3 \\
  Gruber & 0.85 & {--} & {--} & 1.30 & {--} & 1.05 & 3 \\
  Haas & {--} & {--} & {--} & 0.90 & {--} & 1.35 & 2 \\
  Ivanov & 1.00 & 0.95 & 1.15 & {--} & {--} & {--} & 3 \\
\midrule
Mitarbeiter je Qualifikation & 5 & 4 & 3 & 4 & 2 & 3 & \\
\bottomrule
\end{tabular}

  \caption[Qualifikationsmatrix der Fertigung]{Qualifikationsmatrix: Effizienz je Mitarbeiter und Qualifikation. Der Wert 1,00 entspricht der Normalleistung; ein Strich bedeutet, dass die Qualifikation nicht vorliegt und die Zuordnung damit unzulässig ist.}
  \label{tab:qualifikationsmatrix}
\end{table}
\end{landscape}

Tabelle~\ref{tab:qualifikationsmatrix} zeigt die Matrix für die
\ZahlMitarbeiter{} Mitarbeitenden, die in dieser Arbeit betrachtet werden. Zwei
Beobachtungen sind für das weitere Vorgehen wichtig. Erstens ist die Matrix
dünn besetzt: Von den $\ZahlMitarbeiter \times \ZahlQualifikationen = 48$
denkbaren Kombinationen sind nur 21 belegt. Zweitens ist die Besetzung
ungleichmäßig. Fünf Personen dürfen drehen, aber nur zwei sind für die
CNC-Programmierung qualifiziert, und beim Schweißen wie bei der
Qualitätsprüfung sind es jeweils drei. Damit ist bereits an der Matrix
ablesbar, wo Engpässe zu erwarten sind: Arbeitsgänge, die eine seltene
Qualifikation verlangen, konkurrieren um dieselben wenigen Personen,
unabhängig davon, wie gut die übrige Belegschaft ausgelastet ist.

\section{Das Fertigungslos einer Kalenderwoche}
\label{sec:fertigungslos}

Als Datengrundlage der Untersuchung dient das Fertigungslos einer
Kalenderwoche, wie es die \Index{Arbeitsvorbereitung} zu Wochenbeginn freigibt.
Es umfasst \ZahlAufgaben{} Arbeitsgänge, die aus Kundenaufträgen der drei
genannten Branchen hervorgehen. Jeder Arbeitsgang ist durch drei Angaben
bestimmt: eine Bezeichnung, die das zu bearbeitende Teil und die Tätigkeit
benennt, die geforderte Qualifikation und eine \Index{Grundzeit} in Minuten.

Die Grundzeit stammt aus dem Arbeitsplan des jeweiligen Teils. Sie ist die
kalkulierte Bearbeitungszeit bei Normalleistung und schließt Nebenzeiten wie
Aufspannen und Messen ein, nicht jedoch die \Index{Rüstzeit} der Maschine, die im
Unternehmen gesondert erfasst und in dieser Arbeit nicht modelliert wird. Die
Grundzeiten des betrachteten Loses liegen zwischen 45 und 150 Minuten. Diese
Spanne ist typisch für den Betrieb: Die kurzen Arbeitsgänge entfallen
überwiegend auf die Qualitätsprüfung, die langen auf Fräs- und Schweißarbeiten
an größeren Bauteilen.

Aus der Kombination von Fertigungslos und Qualifikationsmatrix ergibt sich die
Größe der Probleminstanz. Von den $\ZahlMitarbeiter \times \ZahlAufgaben = 192$
denkbaren Zuordnungen sind \ZahlPaare{} zulässig; mehr als die Hälfte scheidet
also allein an der fehlenden Qualifikation aus. Für jedes der \ZahlPaare{}
zulässigen Paare ist eine Bearbeitungszeit definiert, für alle übrigen keine.

% Automatisch erzeugt von Code/tabellen.py -- nicht von Hand ändern.
% Quelle: Code/daten.py und Ergebnisse/*.json

\begin{longtable}{@{}llS[table-format=3.0]lS[table-format=3.1]@{}}
\caption{Arbeitsgänge des Fertigungsloses mit der vom Verfahren gefundenen Zuordnung}\label{tab:arbeitsgaenge}\\
\toprule
{Nr.} & {Bezeichnung} & {\text{Grundzeit}} & {Qualifikation} & {\text{Dauer}} \\
 & & {min} & & {min} \\
\midrule
\endfirsthead
\caption[]{Arbeitsgänge (Fortsetzung)}\\
\toprule
{Nr.} & {Bezeichnung} & {\text{Grundzeit}} & {Qualifikation} & {\text{Dauer}} \\
 & & {min} & & {min} \\
\midrule
\endhead
\midrule
\multicolumn{5}{r@{}}{\footnotesize Fortsetzung nächste Seite}\\
\endfoot
\bottomrule
\endlastfoot
  A00 & Flanschwelle A-102 drehen & 95 & Drehen & 95.0 \\
  A01 & Lagerbock B-204 fraesen & 120 & Fräsen & 114.3 \\
  A02 & Rahmen C-310 schweissen & 140 & Schweißen & 121.7 \\
  A03 & Passfläche D-011 schleifen & 60 & Schleifen & 66.7 \\
  A04 & Programm E-556 erstellen & 110 & CNC-Programmierung & 78.6 \\
  A05 & Charge F-078 pruefen & 45 & Qualitätsprüfung & 33.3 \\
  A06 & Antriebswelle A-115 drehen & 130 & Drehen & 118.2 \\
  A07 & Getriebedeckel B-221 fraesen & 85 & Fräsen & 65.4 \\
  A08 & Traverse C-333 schweissen & 105 & Schweißen & 77.8 \\
  A09 & Dichtsitz D-042 schleifen & 75 & Schleifen & 83.3 \\
  A10 & Programm E-560 erstellen & 90 & CNC-Programmierung & 64.3 \\
  A11 & Charge F-081 pruefen & 55 & Qualitätsprüfung & 40.7 \\
  A12 & Spindel A-127 drehen & 70 & Drehen & 56.0 \\
  A13 & Konsole B-238 fraesen & 150 & Fräsen & 115.4 \\
  A14 & Behaelter C-350 schweissen & 95 & Schweißen & 70.4 \\
  A15 & Laufflaeche D-063 schleifen & 110 & Schleifen & 84.6 \\
  A16 & Programm E-571 erstellen & 65 & CNC-Programmierung & 59.1 \\
  A17 & Charge F-090 pruefen & 80 & Qualitätsprüfung & 69.6 \\
  A18 & Buchse A-134 drehen & 50 & Drehen & 45.5 \\
  A19 & Adapterplatte B-245 fraesen & 70 & Fräsen & 66.7 \\
  A20 & Konsolarm C-361 schweissen & 125 & Schweißen & 92.6 \\
  A21 & Fuehrung D-077 schleifen & 90 & Schleifen & 69.2 \\
  A22 & Programm E-583 erstellen & 100 & CNC-Programmierung & 71.4 \\
  A23 & Charge F-095 pruefen & 60 & Qualitätsprüfung & 57.1 \\
\end{longtable}


Tabelle~\ref{tab:arbeitsgaenge} führt die \ZahlAufgaben{} Arbeitsgänge mit
Bezeichnung, Grundzeit, geforderter Qualifikation und der Dauer auf, die sich
in der später berechneten Zuordnung tatsächlich ergibt. Der Vergleich der
beiden Zeitspalten macht die Wirkung der Effizienz sichtbar: Wo die Dauer unter
der Grundzeit liegt, hat das Verfahren den Arbeitsgang einer Person mit
überdurchschnittlicher Routine zugewiesen; wo sie darüber liegt, war eine
schnellere Person entweder nicht qualifiziert oder bereits stärker belastet als
der Gesamtplan es zulässt. Dass das Verfahren einzelne Arbeitsgänge bewusst
langsamer ausführen lässt, als es möglich wäre, ist kein Fehler, sondern
Ausdruck der Zielsetzung: Minimiert wird nicht die Summe der Bearbeitungszeiten,
sondern die Belastung der am stärksten ausgelasteten Person.

\section{Bisherige Disposition und ihre Schwächen}
\label{sec:disposition}

Die Zuteilung der Arbeitsgänge erfolgt im Unternehmen bislang ohne
rechnergestützte Optimierung. Die Fertigungsleitung öffnet zu Wochenbeginn eine
Tabellenkalkulation, in der die freigegebenen Arbeitsgänge und die verfügbaren
Personen aufgelistet sind, und trägt die Zuordnung von Hand ein. Die
Qualifikationsmatrix existiert dabei durchaus, sie liegt jedoch als eigenes
Tabellenblatt vor und wird beim Eintragen aus dem Gedächtnis herangezogen. Eine
Prüfung, ob die eingetragene Person die geforderte Qualifikation überhaupt
besitzt, findet nicht automatisch statt.

Das Vorgehen folgt einer erkennbaren Regel, auch wenn diese nirgends
niedergeschrieben ist. Zuerst werden die Arbeitsgänge mit seltenen
Qualifikationen verteilt, weil dort kaum Wahlmöglichkeiten bestehen. Danach
werden die verbleibenden Arbeitsgänge der Reihe nach vergeben, wobei die
Fertigungsleitung jeweils die Person auswählt, die den Arbeitsgang am
schnellsten erledigen kann und noch als frei gilt. Diese Regel ist plausibel
und liefert zulässige Pläne. Sie ist zugleich eine \Index{Greedy-Regel} im Sinne
von Kapitel~2: Sie trifft jede Entscheidung lokal und revidiert sie nicht,
wenn sich später herausstellt, dass eine andere Wahl den Gesamtplan verbessert
hätte.

Daraus folgt die zentrale Schwäche. Wer in einem Feld die höchste Effizienz
besitzt, erhält dort tendenziell jeden Arbeitsgang und ist am Ende der Woche
überlastet, während andere Kapazität ungenutzt bleibt. Da das Los erst
abgeschlossen ist, wenn der letzte Arbeitsgang beendet wurde, bestimmt allein
die am stärksten belastete Person den \Index{Makespan}. Eine Zuordnung, die einen
schnellen Mitarbeiter entlastet und dafür einen Arbeitsgang bewusst einer
langsameren Person überträgt, kann den Gesamtplan folglich verkürzen, obwohl
sie die Summe der Bearbeitungszeiten erhöht. Genau diese Möglichkeit ist einer
Regel, die Arbeitsgang für Arbeitsgang die jeweils schnellste Person wählt,
strukturell verschlossen.

Wie groß der dadurch verschenkte Spielraum ist, lässt sich beziffern. Die
nachgebildete Greedy-Regel erreicht auf dem betrachteten Los einen Makespan von
\MakespanGreedy{} Minuten, das exakte Verfahren dagegen \Optimum{} Minuten. Die
Differenz von rund einer halben Stunde entsteht nicht durch zusätzliche
Kapazität, sondern allein durch eine andere Verteilung derselben Arbeit auf
dieselben Personen. Hinzu kommen zwei Nachteile, die sich nicht in Minuten
ausdrücken lassen: Die Zuteilung ist an das Erfahrungswissen einzelner Personen
gebunden und bei deren Abwesenheit nur eingeschränkt reproduzierbar, und sie
ist gegenüber den Beschäftigten nicht begründbar, weil kein nachvollziehbares
Kriterium vorliegt, nach dem sie zustande gekommen ist. Ein Optimierungsmodell
beseitigt beide Nachteile als Nebeneffekt: Es macht die Regel explizit, nach
der geplant wird.
